%%%
%%% Radical "⼍"
%%%
\section*{Radical 14: ``⼍''}\addcontentsline{toc}{section}{Radical 14: ⼍}\addcontentsline{loh}{figure}{\#\#\#\# 14: ⼍}

%%%%%%%%%% 冗 %%%%%%%%%%
\subsection*{冗}\addcontentsline{loh}{figure}{冗}

\begin{Entry}{冗}{4}{⼍}
  \begin{Phonetics}{冗}{rong3}
    \definition{adj.}{supérfluo; redundante | cheio de detalhes triviais | incapaz}
    \definition{s.}{ocupação; negócio}
  \end{Phonetics}
\end{Entry}

\begin{Entry}{冗长}{4,4}{⼍,⾧}
  \begin{Phonetics}{冗长}{rong3chang2}[][HSK 7-9]
    \definition{adj.}{tediosamente longo; extenso; prolixo; extenso | incômodo}
  \end{Phonetics}
\end{Entry}

%%%%%%%%%% 写 %%%%%%%%%%
\subsection*{写}\addcontentsline{loh}{figure}{写}

\begin{Entry}{写}{5}{⼍}
  \begin{Phonetics}{写}{xie3}[][HSK 1]
    \definition{v.}{escrever | compor; escrever (como autor, repórter, etc.) | descrever; retratar | pintar; desenhar | expressar a imagem das coisas através da linguagem e da escrita | desenhar (pintura)}
  \synonymref{编写}{bian1xie3}
  \synonymref{创作}{chuang4zuo4}
  \synonymref{画}{hua4}
  \synonymref{书写}{shu1xie3}
  \synonymref{写作}{xie3zuo4}
  \end{Phonetics}
\end{Entry}

\begin{Entry}{写字}{5,6}{⼍,⼦}
  \begin{Phonetics}{写字}{xie3zi4}
    \definition{v.}{escrever (à mão); praticar caligrafia}
  \end{Phonetics}
\end{Entry}

\begin{Entry}{写字台}{5,6,5}{⼍,⼦,⼝}
  \begin{Phonetics}{写字台}{xie3zi4tai2}[][HSK 6]
    \definition[个,张]{s.}{escrivaninha; secretária; escrivaninha de escrever; uma mesa retangular usada para escrever e trabalhar, com gavetas e algumas com pequenos armários}
  \end{Phonetics}
\end{Entry}

\begin{Entry}{写字匠}{5,6,6}{⼍,⼦,⼕}
  \begin{Phonetics}{写字匠}{xie3zi4 jiang4}
    \definition{s.}{escriba; escritor; calígrafo}
  \end{Phonetics}
\end{Entry}

\begin{Entry}{写字楼}{5,6,13}{⼍,⼦,⽊}
  \begin{Phonetics}{写字楼}{xie3zi4lou2}[][HSK 6]
    \definition{s.}{prédio de escritórios}
  \end{Phonetics}
\end{Entry}

\begin{Entry}{写作}{5,7}{⼍,⼈}
  \begin{Phonetics}{写作}{xie3zuo4}[][HSK 3]
    \definition{v.}{escrever artigos; escrever livros, etc.; também se refere especificamente à criação de obras literárias}
  \end{Phonetics}
\end{Entry}

\begin{Entry}{写真}{5,10}{⼍,⼗}
  \begin{Phonetics}{写真}{xie3zhen1}
    \definition{s.}{retrato; pintura de retrato | representação fiel à realidade; representação fiel}
    \definition{v.}{retratar uma pessoa; desenhar um retrato | descrever algo como é}
  \end{Phonetics}
\end{Entry}

\begin{Entry}{写意}{5,13}{⼍,⼼}
  \begin{Phonetics}{写意}{xie3yi4}
    \definition{s.}{estilo de pintura chinesa com pinceladas livres, caracterizado por expressões vívidas e contornos marcantes}
    \definition{v.}{sugerir (em vez de descrever em detalhes)}
  \end{Phonetics}
  \begin{Phonetics}{写意}{xie4yi4}
    \definition{adj.}{confortável; agradável}
  \end{Phonetics}
\end{Entry}

\begin{Entry}{写照}{5,13}{⼍,⽕}
  \begin{Phonetics}{写照}{xie3zhao4}[][HSK 7-9]
    \definition{s.}{retrato}
    \definition{v.}{retratar (uma pessoa ou personagem)}
  \synonymref{缩影}{suo1ying3}
  \synonymref{写真}{xie3zhen1}
  \end{Phonetics}
\end{Entry}

%%%%%%%%%% 军 %%%%%%%%%%
\subsection*{军}\addcontentsline{loh}{figure}{军}

\begin{Entry}{军}{6}{⼍}
  \begin{Phonetics}{军}{jun1}
    \definition*{s.}{Sobrenome: Jun}
    \definition{s.}{forças armadas; exército; tropas | exército; contingente; muitas pessoas participando de uma atividade | exército; unidades militares}
  \end{Phonetics}
\end{Entry}

\begin{Entry}{军人}{6,2}{⼍,⼈}
  \begin{Phonetics}{军人}{jun1ren2}[][HSK 5]
    \definition[名,位,个]{s.}{soldado; militar; pessoal militar; pessoas com status militar; pessoas servindo nas forças armadas}
  \end{Phonetics}
\end{Entry}

\begin{Entry}{军队}{6,4}{⼍,⾩}
  \begin{Phonetics}{军队}{jun1dui4}[][HSK 6]
    \definition[支,个]{s.}{forças armadas; exército; tropas}
  \end{Phonetics}
\end{Entry}

\begin{Entry}{军事}{6,8}{⼍,⼅}
  \begin{Phonetics}{军事}{jun1shi4}[][HSK 6]
    \definition{s.}{militar; assuntos militares; assuntos relativos aos militares e à guerra}
  \end{Phonetics}
\end{Entry}

\begin{Entry}{军官}{6,8}{⼍,⼧}
  \begin{Phonetics}{军官}{jun1guan1}[][HSK 7-9]
    \definition{s.}{oficial; militares com patente igual ou superior a tenente, também se refere a oficiais com patente igual ou superior a comandante de pelotão nas forças armadas}
  \end{Phonetics}
\end{Entry}

\begin{Entry}{军舰}{6,10}{⼍,⾈}
  \begin{Phonetics}{军舰}{jun1jian4}[][HSK 6]
    \definition[艘,只]{s.}{navio de guerra; embarcação naval | \emph{warcraft}; um termo geral para embarcações militares equipadas com armas e equipamentos que podem executar missões de combate, incluindo principalmente navios de guerra, cruzadores, contratorpedeiros, porta"-aviões, submarinos, torpedeiros, etc.}
  \end{Phonetics}
\end{Entry}

\begin{Entry}{军装}{6,12}{⼍,⾐}
  \begin{Phonetics}{军装}{jun1zhuang1}
    \definition[套,件,身]{s.}{uniforme militar (ou do exército); uniforme}
  \synonymref{礼服}{li3fu2}
  \end{Phonetics}
\end{Entry}

%%%%%%%%%% 农 %%%%%%%%%%
\subsection*{农}\addcontentsline{loh}{figure}{农}

\begin{Entry}{农}{6}{⼍}
  \begin{Phonetics}{农}{nong2}
    \definition*{s.}{Sobrenome: Nong}
    \definition{s.}{agricultura; criação de animais | camponês; fazendeiro}
  \end{Phonetics}
\end{Entry}

\begin{Entry}{农历}{6,4}{⼍,⼚}
  \begin{Phonetics}{农历}{nong2li4}[][HSK 7-9]
    \definition{s.}{o calendário lunar; o calendário tradicional chinês}
  \end{Phonetics}
\end{Entry}

\begin{Entry}{农业}{6,5}{⼍,⼀}
  \begin{Phonetics}{农业}{nong2ye4}[][HSK 3]
    \definition{s.}{agricultura}
  \end{Phonetics}
\end{Entry}

\begin{Entry}{农民}{6,5}{⼍,⽒}
  \begin{Phonetics}{农民}{nong2min2}[][HSK 3]
    \definition[个,位,名,些]{s.}{fazendeiro; camponês; campesinato; trabalhadores que participam da produção agrícola há muito tempo}
  \end{Phonetics}
\end{Entry}

\begin{Entry}{农民工}{6,5,3}{⼍,⽒,⼯}
  \begin{Phonetics}{农民工}{nong2min2gong1}[][HSK 7-9]
    \definition{s.}{trabalhadores migrantes; trabalhadores com registro de domicílio agrícola que exercem atividades não agrícolas em áreas urbanas}
  \end{Phonetics}
\end{Entry}

\begin{Entry}{农产品}{6,6,9}{⼍,⼇,⼝}
  \begin{Phonetics}{农产品}{nong2chan3pin3}[][HSK 5]
    \definition[批]{s.}{produtos agrícolas}
  \end{Phonetics}
\end{Entry}

\begin{Entry}{农场}{6,6}{⼍,⼟}
  \begin{Phonetics}{农场}{nong2chang3}[][HSK 7-9]
    \definition[座,家]{s.}{fazenda; empresas que utilizam máquinas e se dedicam à produção agrícola em larga escala}
  \end{Phonetics}
\end{Entry}

\begin{Entry}{农作物}{6,7,8}{⼍,⼈,⽜}
  \begin{Phonetics}{农作物}{nong2zuo4wu4}[][HSK 7-9]
    \definition{s.}{culturas agrícolas; plantas agrícolas; termo genérico para diversas culturas agrícolas, como grãos, oleaginosas, hortaliças, algodão e tabaco}
  \end{Phonetics}
\end{Entry}

\begin{Entry}{农村}{6,7}{⼍,⽊}
  \begin{Phonetics}{农村}{nong2cun1}[][HSK 3]
    \definition{s.}{aldeia; campo; área rural; locais onde vivem os trabalhadores principalmente dedicados à produção agrícola}
  \end{Phonetics}
\end{Entry}

%%%%%%%%%% 冠 %%%%%%%%%%
\subsection*{冠}\addcontentsline{loh}{figure}{冠}

\begin{Entry}{冠}{9}{⼍}
  \begin{Phonetics}{冠}{guan1}
    \definition{s.}{chapéu | corona; coroa; copa | crista}
  \end{Phonetics}
  \begin{Phonetics}{冠}{guan4}
    \definition*{s.}{Sobrenome: Guan}
    \definition{s.}{primeiro lugar; o melhor; classificado em primeiro lugar}
    \definition{v.}{colocar um chapéu (boné) | preceder com (por); coroar com; adicionar um nome ou texto na frente}
  \end{Phonetics}
\end{Entry}

\begin{Entry}{冠军}{9,6}{⼍,⼍}
  \begin{Phonetics}{冠军}{guan4jun1}[][HSK 5]
    \definition[位,名,项,个]{s.}{campeão; medalhista de ouro; primeiro lugar em esportes e outras competições}
  \end{Phonetics}
\end{Entry}

%%%%%%%%%% 冤 %%%%%%%%%%
\subsection*{冤}\addcontentsline{loh}{figure}{冤}

\begin{Entry}{冤}{10}{⼍}
  \begin{Phonetics}{冤}{yuan1}[][HSK 7-9]
    \definition{adj.}{que não vale a pena; inútil; um desperdício (de tempo/dinheiro/esforço) ; custou tempo e dinheiro, mas não atingiu o objetivo; não valeu a pena | sentir"-se injustiçado; sentir ressentimento; isso descreve a sensação de ser criticado ou punido por algo que você não fez; sentir que é injusto}
    \definition{s.}{tratamento injusto; acusação injusta; injustiça; erro judiciário; condenação injusta | rancor; ressentimento; ódio; rixa; inimizade}
    \definition{v.}{enganar; ludibriar; iludir | ser falsamente acusado; ser injustiçado; ser tratado injustamente ou causar injustiça a alguém}
  \end{Phonetics}
\end{Entry}

\begin{Entry}{冤枉}{10,8}{⼍,⽊}
  \begin{Phonetics}{冤枉}{yuan1wang5}[][HSK 7-9]
    \definition{adj.}{injusto; injustificado; descreve a sensação de ser criticado injustamente quando você não fez nada de errado | indigno; não merecido; descreve o uso de tempo, dinheiro, etc., como algo que não vale a pena}
    \definition{v.}{prejudicar; tratar injustamente; ser acusado falsamente}
  \synonymref{勉强}{mian3qiang3}
  \synonymref{曲折}{qu1zhe2}
  \synonymref{委屈}{wei3qu5}
  \synonymref{无辜}{wu2gu1}
  \antonymref{美化}{mei3hua4}
  \end{Phonetics}
\end{Entry}

%%%%% EOF %%%%%

